\documentclass{ohm_project_description}

\projecttitle{M-APR: Emergente Systeme in der Produktion}
\projectauthor{TTZ Nürnberger Land / TH Nürnberg}
\projectdate{2026}

\begin{document}

\maketitle
\thispagestyle{fancy}

Moderne Produktionsumgebungen setzen zunehmend auf heterogene Roboterflotten — mobile Transportplattformen (AMR), kollaborative Roboterarme und stationäre Bearbeitungsstationen —, die gemeinsam Fertigungs- und Logistikaufgaben bewältigen. Die Koordination erfolgt heute überwiegend zentral. Bei Störungen oder Maschinenausfällen ist eine zentrale Koordination jedoch anfällig: Fällt der zentrale Knoten aus, steht die gesamte Prozesskette still.

Dieses M-APR-Projekt (3 Semester, Start WiSe 2026/2027) untersucht, unter welchen Bedingungen dezentrale, emergente Koordinationsstrategien einer zentralen Steuerung in Produktionsszenarien überlegen sind. Im Fokus steht der systematische Vergleich zentraler und dezentraler Ansätze unter kontrolliert variierten Störbedingungen (Kommunikationsausfälle, Maschinenausfälle, dynamische Auftragslage). Im dritten Semester werden ausgewählte Szenarien auf reale Hardware am TTZ Nürnberger Land übertragen.

\section*{Arbeitspakete (3 Semester)}
\begin{itemize}[leftmargin=0.5cm]
    \setlength\itemsep{.1em}
    \item \textbf{Semester 1:} Literaturrecherche, Aufbau der ROS-2-Simulationsumgebung, zentrale Baseline und erstes dezentrales Verfahren
    \item \textbf{Semester 2:} Multi-Agent Reinforcement Learning (MARL), systematischer Vergleich unter Störbedingungen, statistische Auswertung
    \item \textbf{Semester 3:} Transfer auf reale Hardware am TTZ, Sim-to-Real-Evaluation, Masterarbeit und Konferenzpublikation
\end{itemize}

\section*{Voraussetzungen}
\begin{itemize}[leftmargin=0.5cm]
    \setlength\itemsep{.1em}
    \item Bachelorabschluss in EI, Informatik, Mechatronik, Robotik o.\,ä.
    \item Programmierkenntnisse in C++ und/oder Python
    \item Grundkenntnisse in Machine Learning / Reinforcement Learning
    \item Erfahrung mit ROS\,2 und Simulation (Gazebo o.\,ä.) von Vorteil
\end{itemize}

\vspace{0.5cm}
Das Thema ist als \textbf{M-APR (3 Semester)} konzipiert, kann nach Abstimmung auch als Projekt- oder Masterarbeit bearbeitet werden.

\vfill
\textcolor{ohm_red}{\rule{\linewidth}{0.4mm}}
\textbf{\textcolor{ohm_red}{TTZ Nürnberger Land / Labor für mobile Robotik}} \\
\begin{tabular}{@{}ll}
\textbf{Betreuer:}   & Prof. Dr. Christian Pfitzner \\
\textbf{E-Mail:}     & \href{mailto:christian.pfitzner@th-nuernberg.de}{christian.pfitzner@th-nuernberg.de} \\
\textbf{Co-Betreuer:} & Waldemar Haag \\
\textbf{E-Mail:}     & \href{mailto:waldemar.haag@th-nuernberg.de}{waldemar.haag@th-nuernberg.de} \\
\end{tabular}

\end{document}
